\documentclass[11pt,a4paper]{article}
\usepackage[margin=1in]{geometry}
\usepackage{amsmath}
\usepackage{amssymb}
\usepackage{booktabs}
\usepackage{hyperref}
\usepackage{graphicx}
\usepackage{xcolor}
\usepackage{setspace}

\setlength{\parskip}{0.5em}
\setlength{\parindent}{0pt}
\onehalfspacing

\title{Computer Animation HW2 Report\\Skeleton Animation and Interpolation}
\author{Student ID: 112550052}
\date{\today}

\begin{document}
\maketitle

\section{Introduction and Motivation}
In this homework, I implemented a full skeleton animation pipeline based on quaternion rotations.
The main goal was to replace Euler-angle based joint rotation handling with quaternion-based processing, then build interpolation on top of it.
Compared with Euler angles, quaternions are more stable for continuous animation and avoid gimbal lock.

I focused on five required parts and one bonus extension:
\begin{itemize}
    \item converting AMC Euler rotations to unit quaternions,
    \item forward kinematics in a hierarchy tree,
    \item keyframe hemisphere normalization,
    \item root position interpolation in world space,
    \item quaternion interpolation with Slerp and Squad,
    \item bonus: root interpolation method toggle (Linear vs Catmull-Bezier).
\end{itemize}

\section{Fundamental Implementation}

\subsection{Part 1: Euler to Quaternion}
In \texttt{motion.cpp}, each bone rotation read from AMC is converted from Euler degrees to quaternion.
I used the same axis convention as the provided helper utilities, then normalized the quaternion before storing.
This step is important because all downstream interpolation and FK assume valid unit quaternions.

\subsection{Part 2: Forward Kinematics}
In \texttt{kinematics.cpp}, I implemented recursive traversal over the skeleton hierarchy:
\begin{itemize}
    \item root bone starts from frame root translation,
    \item each child starts at its parent end position,
    \item global rotation is accumulated from parent global rotation and local joint rotation,
    \item end position is computed from rotated local direction times bone length.
\end{itemize}
This gives consistent global start/end positions and global orientation for every bone in each frame.

\subsection{Part 3: Hemisphere Normalization}
In \texttt{main.cpp} keyframe preprocessing, I added a global pass over keyframes.
For each bone, if the dot product between two consecutive keyframe quaternions is negative, I flip the current quaternion sign.
Then I normalize again.

This keeps neighboring keyframes on the same hemisphere and avoids taking the long arc during interpolation.
Without this step, sudden 180-degree-like flips are likely to appear.

\subsection{Part 4: RootPositionInterpolation}
I implemented root translation interpolation in \texttt{interpolation\_math.cpp} using cubic Bezier with Catmull-Rom style tangents.
Given \((k_{m1}, k_0, k_1, k_2)\), I compute
\[
m_0 = \frac{1}{2}(k_1-k_{m1}), \quad
m_1 = \frac{1}{2}(k_2-k_0)
\]
and convert to Bezier control points:
\[
c_0 = k_0,\quad
c_1 = k_0 + \frac{m_0}{3},\quad
c_2 = k_1 - \frac{m_1}{3},\quad
c_3 = k_1.
\]
Then the interpolated root position is \(B(t)\), the cubic Bezier curve at parameter \(t\in[0,1]\).

This produces smoother root movement than simple linear interpolation, especially around keyframe boundaries.

\subsection{Part 5: Quaternion Interpolation}
\textbf{Slerp:}
I implemented robust Slerp with:
\begin{itemize}
    \item input normalization,
    \item hemisphere fix using dot-product sign,
    \item near-parallel fallback to normalized linear blend (to avoid numerical instability),
    \item standard trigonometric Slerp formula otherwise.
\end{itemize}

\textbf{Squad (basic version):}
I implemented Squad by combining nested Slerps:
\begin{itemize}
    \item one arc between the segment endpoints,
    \item one control arc computed from neighboring keyframes,
    \item final blend using factor \(2t(1-t)\).
\end{itemize}
This gives smoother rotational transition than plain Slerp in many joints.

\section{Questions and Discussion}

\subsection{Q1: Introduce your implementation in RootPositionInterpolation()}
I used cubic Bezier interpolation with Catmull-Rom-derived tangents.
The idea is to keep the segment endpoints fixed at \(k_0\) and \(k_1\), while estimating incoming/outgoing motion direction from neighboring keyframes \(k_{m1}\) and \(k_2\).
Compared with linear interpolation, this method gives better continuity of velocity at keyframe boundaries and visually smoother trajectory.

\subsection{Q2: Observation and comparison between Slerp and Squad}
My observations from interactive playback are:
\begin{itemize}
    \item Slerp is stable and simple, and each segment follows the shortest quaternion arc.
    \item However, at keyframe boundaries, Slerp may still feel slightly ``piecewise'' because each segment is independent.
    \item Squad produces smoother transitions across boundaries because it uses neighboring keyframes to shape the rotation curve.
    \item In clips with larger spacing between keyframes, the visual difference becomes clearer (Squad generally looks more natural).
\end{itemize}

\subsection{Q3: Why do we need hemisphere-normalized keyframe rotations?}
Quaternions \(q\) and \(-q\) represent the same physical rotation, but interpolation depends on the 4D direction in quaternion space.
If two adjacent keyframes use opposite signs, interpolation may travel along a long arc, causing abrupt flips.
Hemisphere normalization enforces sign consistency across the keyframe sequence, so Slerp/Squad follows the expected short path and motion stays continuous.

\subsection{Discussion and Findings}
To make the comparison more visible, I tested the same animation clips with different keyframe intervals (\(10\), \(20\), and \(35\)).
When the interval was \(10\), both Slerp and Squad looked acceptable because the keyframes were already dense.
At interval \(20\), the difference became easier to notice: Slerp still worked, but joint motion around shoulders and hips looked slightly segmented near keyframe transitions.
At interval \(35\), the contrast was clear. Squad preserved smoother rotational flow, while Slerp looked more ``connect-the-dots'' between segments.

For root motion, cubic Bezier interpolation gave a smoother trajectory than linear interpolation, especially when the character changed moving direction.
I also found that hemisphere normalization is not just a theoretical detail: without it, large visual flips can happen even when the keyframe poses themselves are valid.
After normalization, both Slerp and Squad became much more stable and predictable during playback.

\section{Bonus Implementation}
I implemented a bonus feature in the control panel: a root interpolation toggle that switches between \texttt{Linear} and \texttt{Catmull-Bezier}.
This was done by adding an extra root interpolation function (\texttt{RootPositionLinear}) and wiring the UI button to re-bake keyframe data using the selected position interpolator.

In terms of behavior, the linear mode is useful as a baseline because its result is easy to predict.
The Catmull-Bezier mode usually gives smoother turns and less abrupt speed change near keyframe boundaries.
This side-by-side toggle made it much easier to compare interpolation behavior and confirm the effect of higher-order position interpolation.

\section{Conclusion}
This homework helped me understand why quaternion-native pipelines are standard in character animation.
After replacing Euler-based rotation storage and adding hemisphere normalization, interpolation became much more reliable.
Slerp is a good baseline, while Squad improves smoothness at segment boundaries.
For root translation, cubic Bezier with Catmull-style tangents provides a clean and practical trade-off between smoothness and implementation complexity.

\end{document}
